%% toc: P15 158 158.6,158.7
% Proofs for test-function weak convergence on topological spaces and the point-mass boundary remark (Durrett, Section 3.2, pp. 119--120), nice spaces and the extension theorem (Durrett, pp. 53--54), the cardinality of Borel sets (Durrett, p. 54), and the birthday error bound and density notation (Durrett, pp. 117--118).
\section{Test functions, nice spaces and Borel sets: proofs}
\label{sec:annex-p15-w02c-p15-004-weak-test-functions}

Use the definitions of
Chapter~\ref{sec:p15-w02c-p15-004-weak-test-functions} and
\cref{sec:p15-weak-convergence-of-laws}. Two facts are used repeatedly.
First, if \(f:S\to T\) is a map between topological spaces and
\(f^{-1}(V)\) is open for every open \(V\subseteq T\), then \(f\) is
\(\mathcal B(S)/\mathcal B(T)\)-measurable: the open sets of \(T\)
generate \(\mathcal B(T)\), each of their inverse images is open and so
lies in \(\mathcal B(S)\), and \cref{thm:p15-measurability-criterion}
gives measurability. Second, for a point mass,
\(\int g\,d\delta_x=g(x)\) for every bounded measurable \(g\): the set
\(\{y:g(y)\neq g(x)\}\) does not contain \(x\) and so is
\(\delta_x\)-null, hence \(g\) equals the constant \(g(x)\)
\(\delta_x\)-almost everywhere, and a constant integrates to itself
against a probability measure.

\begin{annexitem}{annex:p15-w02c-p15-004-topo}{thm:p15-w02c-p15-004-topological-weak}{Meaning, the line, images and uniqueness}
\begin{proof}
(i) A continuous \(g:S\to\mathbb R\) pulls open subsets of
\(\mathbb R\) back to open subsets of \(S\), so it is Borel measurable
by the first preliminary fact, the open sets of \(\mathbb R\) generating
\(\mathcal B(\mathbb R)\) by \cref{def:p15-borel-sets}. If also
\(|g|\leq M<\infty\), then \(g^+\leq M\) and \(g^-\leq M\), so each has
integral at most \(M\mu(S)=M\) by monotonicity of the nonnegative
integral. Thus \(g\) is integrable, and the absolute-value bound of
\cref{thm:p15-signed-integral-laws} gives
\(|\int g\,d\mu|\leq\int|g|\,d\mu\leq M\). Each integral in the
definition is therefore a finite real number, and the convergence
requirement is a statement about ordinary real sequences.
If \(X_n\) is an \(S\)-valued random element on a probability space
\((\Omega_n,\mathcal F_n,\mathbb P_n)\) with law \(\mu_n\), then
\(g(X_n)=g\circ X_n\) is a real random variable by composition
(\cref{thm:p15-measurable-operations}), it is bounded, and
\cref{thm:p15-expectation-law-integral} gives
\(\mathbb E g(X_n)=\int_S g\,d\mu_n\). The probability spaces may differ
with \(n\); only the laws enter.

(ii) For \(S=\mathbb R\) with the usual topology, the continuous
functions of \cref{def:p15-w02c-p15-004-topological-weak} are those
whose inverse images of open sets are open, which is the usual
\(\varepsilon\)--\(\delta\) continuity on \(\mathbb R\): if inverse
images of open sets are open, the inverse image of
\((g(x)-\varepsilon,g(x)+\varepsilon)\) is an open set containing
\(x\) and so contains some interval \((x-\delta,x+\delta)\); conversely,
under \(\varepsilon\)--\(\delta\) continuity each point of
\(g^{-1}(U)\) has such an interval inside \(g^{-1}(U)\). The Borel
sigma-algebra is the one of \cref{def:p15-borel-sets}. The condition of
the definition is therefore literally the second condition of
\cref{thm:p15-weak-bounded-continuous}, which that theorem proves
equivalent to convergence of distribution functions at every continuity
point of \(F_\infty\).

(iii) If \(V\subseteq T\) is open then \(h^{-1}(V)\) is open, so
\(h\) is \(\mathcal B(S)/\mathcal B(T)\)-measurable by the first
preliminary fact. By \cref{thm:p15-pushforward-laws}, applied on the
probability space \((S,\mathcal B(S),\mu_n)\) to the random element
\(h\), each \(\nu_n\defeq\mu_n\circ h^{-1}\) is a Borel probability
measure on \(T\). Let \(g\in C_b(T)\). For open \(U\subseteq\mathbb R\),
\((g\circ h)^{-1}(U)=h^{-1}(g^{-1}(U))\) is open, and
\(\sup_S|g\circ h|\leq\sup_T|g|\); hence \(g\circ h\in C_b(S)\).
\Cref{thm:p15-expectation-law-integral}, on the same probability space
and random element, gives \(\int_T g\,d\nu_n=\int_S g\circ h\,d\mu_n\)
for every \(n\), including \(n=\infty\). Therefore
\[
 \int_T g\,d\nu_n=\int_S g\circ h\,d\mu_n
   \longrightarrow\int_S g\circ h\,d\mu_\infty=\int_T g\,d\nu_\infty,
\]
the convergence being the hypothesis applied to \(g\circ h\). As \(g\)
was arbitrary, \(\nu_n\Rightarrow\nu_\infty\).

(iv) Let \(d\) induce the topology: a set is open exactly when each of
its points is the centre of some open ball \(\{y:d(x,y)<r\}\), \(r>0\),
contained in the set. Fix a nonempty closed set \(K\) and put
\(d(x,K)\defeq\inf_{y\in K}d(x,y)\). For \(x,z\in S\) and \(y\in K\),
\(d(x,y)\leq d(x,z)+d(z,y)\); taking the infimum over \(y\) gives
\(d(x,K)\leq d(x,z)+d(z,K)\), and by symmetry
\(|d(x,K)-d(z,K)|\leq d(x,z)\). Hence \(x\mapsto d(x,K)\) is continuous:
the inverse image of an open \(U\subseteq\mathbb R\) contains, with each
point \(x\), the ball of radius \(r\) about \(x\) whenever
\((d(x,K)-r,d(x,K)+r)\subseteq U\). If \(x\in K\), then
\(d(x,K)=0\). If \(x\notin K\), the open set \(K^c\) contains a ball of
some radius \(r>0\) about \(x\), so \(d(x,y)\geq r\) for every
\(y\in K\) and \(d(x,K)\geq r>0\). For \(k=1,2,\ldots\) put
\[
 f_k(x)=\max\{0,\,1-k\,d(x,K)\}.
\]
Each \(f_k\) is a composition of the continuous function \(d(\cdot,K)\)
with the continuous real map \(t\mapsto\max\{0,1-kt\}\), so it is
continuous, and \(0\leq f_k\leq1\); thus \(f_k\in C_b(S)\). If
\(x\in K\), then \(f_k(x)=1\) for every \(k\). If \(x\notin K\), then
\(f_k(x)=0\) as soon as \(k\geq1/d(x,K)\). So \(f_k\to\mathbf1_K\)
pointwise, with \(|f_k|\leq1\).

Now suppose \(\mu_n\Rightarrow\mu\) and \(\mu_n\Rightarrow\nu\). For
each \(k\), the real sequence \(\int f_k\,d\mu_n\) converges both to
\(\int f_k\,d\mu\) and to \(\int f_k\,d\nu\), so these are equal. The
bounded convergence theorem (\cref{thm:p15-dominated-bounded-convergence},
first part with the constant dominating function \(1\), integrable
because the measures are finite) gives
\[
 \mu(K)=\lim_k\int f_k\,d\mu=\lim_k\int f_k\,d\nu=\nu(K).
\]
For \(K=\varnothing\) both sides are zero. So \(\mu\) and \(\nu\) agree
on the collection \(\mathcal K\) of closed sets. This collection is a
pi-system: the intersection of two closed sets has as complement the
union of two open sets, which is open. It generates \(\mathcal B(S)\):
\(\sigma(\mathcal K)\) contains the complement of each closed set, that
is, every open set, so \(\mathcal B(S)\subseteq\sigma(\mathcal K)\);
and every closed set is the complement of an open set, so
\(\sigma(\mathcal K)\subseteq\mathcal B(S)\).
\Cref{thm:p15-measure-uniqueness} now gives \(\mu=\nu\) on
\(\mathcal B(S)\).
\end{proof}
\end{annexitem}

\begin{annexitem}{annex:p15-w02c-p15-004-nonunique}{thm:p15-w02c-p15-004-nonunique}{The two-point counterexample}
\begin{proof}
The collection \(\mathcal T=\{\varnothing,\{a\},S\}\) contains
\(\varnothing\) and \(S\). The intersection of any two members is the
smaller of them, and the union of any family of members is the largest
member in the family (or \(\varnothing\) for the empty family); since
the members are totally ordered by inclusion, both lie in
\(\mathcal T\). So \(\mathcal T\) is a topology. The sigma-algebra
\(\sigma(\mathcal T)\) contains \(\{a\}\) and its complement \(\{b\}\),
hence all four subsets of \(S\).

Let \(g:S\to\mathbb R\) be continuous and suppose \(g(a)\neq g(b)\).
With \(r=|g(a)-g(b)|>0\), the open interval
\(U=(g(b)-r,g(b)+r)\) contains \(g(b)\) but not \(g(a)\), so
\(g^{-1}(U)=\{b\}\), which is not in \(\mathcal T\). This contradicts
continuity, so \(g(a)=g(b)\) and \(g\) is constant. Conversely each
constant function \(c\) is continuous, since \(c^{-1}(U)\) is
\(\varnothing\) or \(S\), and it is bounded. Hence \(C_b(S)\) is exactly
the set of constants.

For a constant \(c\) and any Borel probability measure \(\mu\) on \(S\),
\(\int c\,d\mu=c\,\mu(S)=c\). For any sequence \(\mu_n\) and any
\(\mu_\infty\), the sequences \(\int g\,d\mu_n\) are therefore the
constant sequence \(c\), equal to \(\int g\,d\mu_\infty\), for every
\(g\in C_b(S)\). This proves that every sequence converges to every
Borel probability measure, and in particular that the constant sequence
\(\delta_a\) converges to \(\delta_a\) and to \(\delta_b\). These two
measures are distinct because \(\delta_a(\{a\})=1\) and
\(\delta_b(\{a\})=0\).

Finally, if a metric \(d\) induced \(\mathcal T\), then part~(iv) of
\cref{thm:p15-w02c-p15-004-topological-weak} would force
\(\delta_a=\delta_b\), which is false. The same conclusion can be seen
directly: in a metric space the ball of radius \(d(a,b)>0\) about \(b\)
is an open set containing \(b\) and not \(a\), but the only member of
\(\mathcal T\) containing \(b\) is \(S\).
\end{proof}
\end{annexitem}

\begin{annexitem}{annex:p15-w02c-p15-004-point-masses}{thm:p15-w02c-p15-004-point-masses}{Both directions for point masses}
\begin{proof}
By \cref{thm:p15-weak-bounded-continuous}, \(\delta_{x_n}\Rightarrow
\delta_{x_\infty}\) holds exactly when
\(\int g\,d\delta_{x_n}\to\int g\,d\delta_{x_\infty}\) for every
bounded continuous \(g:\mathbb R\to\mathbb R\), that is, by the second
preliminary fact, exactly when \(g(x_n)\to g(x_\infty)\) for every such
\(g\).

If \(x_n\to x_\infty\), continuity of \(g\) at \(x_\infty\) gives
\(g(x_n)\to g(x_\infty)\) for every bounded continuous \(g\), so the
point masses converge weakly.

Conversely, assume weak convergence and take
\(g(y)=\min\{1,|y-x_\infty|\}\). It is bounded by one and continuous,
being a composition of the continuous maps \(y\mapsto|y-x_\infty|\) and
\(t\mapsto\min\{1,t\}\). Hence
\(\min\{1,|x_n-x_\infty|\}=g(x_n)\to g(x_\infty)=0\). Given
\(0<\varepsilon<1\), for all sufficiently large \(n\) this minimum is
below \(\varepsilon\), and since \(\varepsilon<1\) the minimum is then
\(|x_n-x_\infty|\). So \(x_n\to x_\infty\).

For the random-variable form, \(\mathbb P(X_n=x_n)=1\) says that the
law of \(X_n\) gives mass one to \(\{x_n\}\), hence mass
\(\mathbf1_B(x_n)\) to every Borel \(B\); the law is
\(\delta_{x_n}\), and \(X_n\Rightarrow X_\infty\) means
\(\delta_{x_n}\Rightarrow\delta_{x_\infty}\).
\end{proof}
\end{annexitem}

\begin{annexitem}{annex:p15-w02c-p15-004-boundary-approach}{thm:p15-w02c-p15-004-boundary-approach}{The general boundary approach and its strict inequalities}
\begin{proof}
Here the closure \(\overline A\) of \(A\subseteq\mathbb R\) is the
intersection of all closed sets containing \(A\), and the interior
\(A^\circ\) is the union of all open sets contained in \(A\).

(i) Since \(G\) is open, \(G^\circ=G\), so
\(\partial G=\overline G\setminus G\) and \(x_\infty\notin G\). Let
\(n\geq1\) and \(I_n=(x_\infty-1/n,x_\infty+1/n)\). If \(I_n\cap G\)
were empty, \(I_n^c\) would be a closed set containing \(G\), so
\(\overline G\subseteq I_n^c\) and \(x_\infty\notin\overline G\),
contrary to \(x_\infty\in\partial G\subseteq\overline G\). Choose
\(x_n\in I_n\cap G\). Then \(x_n\in G\) and
\(|x_n-x_\infty|<1/n\), so \(x_n\to x_\infty\).

(ii) Let \(x_n\in G\) be any sequence with \(x_n\to x_\infty\). By
\cref{thm:p15-w02c-p15-004-point-masses}, \(X_n\Rightarrow X_\infty\).
The law of \(X_n\) is \(\delta_{x_n}\), so
\(\mathbb P(X_n\in G)=\mathbf1_G(x_n)=1\), and
\(\mathbb P(X_\infty\in G)=\mathbf1_G(x_\infty)=0\) by (i). The
\(\liminf\) of the constant sequence one is one, which exceeds zero.

(iii) \(K=G^c\) is closed because its complement \(G\) is open.
Since \(x_n\in G\), \(\mathbb P(X_n\in K)=0\) for all \(n\), and since
\(x_\infty\notin G\), \(\mathbb P(X_\infty\in K)=1\). The
\(\limsup\) of the zero sequence is zero, which is less than one.

(iv) \(K\) is closed, so \(\overline K=K=G^c\). A set open and
contained in \(K\) is disjoint from \(G\), so its complement is a
closed set containing \(G\), hence containing \(\overline G\); thus
\(K^\circ\subseteq(\overline G)^c\). Conversely \((\overline G)^c\) is
open and contained in \(G^c=K\). So \(K^\circ=(\overline G)^c\) and
\[
 \partial K=\overline K\setminus K^\circ
          =G^c\cap\overline G=\overline G\setminus G=\partial G.
\]
As \(x_\infty\in\partial G\), \(\mathbb P(X_\infty\in\partial G)=1\neq0\),
so neither set is a continuity set for \(\delta_{x_\infty}\) in the
sense of \cref{def:p15-weak-boundary-sets}. Part~(iv) of
\cref{thm:p15-weak-portmanteau} concerns only continuity sets, and by
(ii) the probabilities of \(G\) are one along the sequence and zero in
the limit, so they do not converge to it.

(v) If \(x_\infty\in G\), openness of \(G\) gives \(r>0\) with
\((x_\infty-r,x_\infty+r)\subseteq G\). Since \(x_n\to x_\infty\),
\(|x_n-x_\infty|<r\) for all sufficiently large \(n\), and then
\(x_n\in G\) and \(\mathbb P(X_n\in G)=1\). Hence
\(\mathbb P(X_n\in G)\to1=\mathbb P(X_\infty\in G)\).
\end{proof}
\end{annexitem}

\begin{annexitem}{annex:p15-w02c-p15-004-nice-coding}{thm:p15-w02c-p15-004-nice-coding}{Nice maps against Borel real codings}
\begin{proof}
Recall from \cref{def:p15-measurable-real-coding} that a Borel real coding
of \((S,\mathcal S)\) onto a Borel set \(E\subseteq\mathbb R\) is a
bijection \(\varphi:S\to E\) that is measurable from \((S,\mathcal S)\) to
\(E\) with its trace sigma-algebra \(\{E\cap B:B\in\mathcal B(\mathbb R)\}\),
and whose inverse is measurable in the opposite direction.

(i) Let \(\varphi\) be such a coding. It is a bijection onto \(E\), so
\(\varphi(S)=E\), and it is injective as a map into \(\mathbb R\). For
\(B\in\mathcal B(\mathbb R)\),
\(\varphi^{-1}(B)=\varphi^{-1}(E\cap B)\in\mathcal S\), because
\(E\cap B\) is a member of the trace sigma-algebra of \(E\). The inverse
is defined on \(\varphi(S)=E\) and is measurable from \(E\) with its trace
sigma-algebra to \((S,\mathcal S)\) by hypothesis. These are the clauses of
\cref{def:p15-w02c-p15-004-nice}.

(ii) Let \(\varphi\) be a nice map. If \(\varphi\) is a Borel real coding
onto \(\varphi(S)\), then \(\varphi(S)\) is Borel by the definition of a
coding. Conversely, let \(\varphi(S)\in\mathcal B(\mathbb R)\). The map
\(\varphi\) is a bijection of \(S\) onto \(\varphi(S)\). For a member
\(\varphi(S)\cap B\) of the trace sigma-algebra,
\(\varphi^{-1}(\varphi(S)\cap B)=\varphi^{-1}(B)\in\mathcal S\), so
\(\varphi\) is measurable into \(\varphi(S)\) with its trace
sigma-algebra; the inverse is measurable by the definition of a nice map.
Hence \(\varphi\) is a Borel real coding onto the Borel set
\(\varphi(S)\).

(iii) Let \(\varphi\) be a nice map. If \(\varphi(S)\) is Borel and
\(A\in\mathcal S\), then \(\varphi(A)\) is the inverse image of \(A\) under
the measurable map \(\varphi^{-1}:\varphi(S)\to S\), so
\(\varphi(A)=\varphi(S)\cap B\) for some \(B\in\mathcal B(\mathbb R)\), an
intersection of two Borel sets. Conversely, if every \(\varphi(A)\) is
Borel, take \(A=S\in\mathcal S\).

(iv) Let \(\iota:A\to\mathbb R\) be the inclusion. It is injective, and
\(\iota^{-1}(B)=A\cap B\) lies in the trace sigma-algebra of \(A\) for
every Borel \(B\). Its image is \(A\), and its inverse is the identity of
\(A\), which is measurable from \(A\) with the trace sigma-algebra to
\(A\) with the same sigma-algebra. So \(\iota\) is a nice map, and by
part~(ii) it is a Borel real coding exactly when \(\iota(A)=A\) is Borel.
\end{proof}
\end{annexitem}

\begin{annexitem}{annex:p15-w02c-p15-004-durrett-2122}{thm:p15-w02c-p15-004-durrett-2122}{Three readings of the Borel sets of a subspace, and the coding}
\begin{proof}
(i) For \(x\in S\) and \(r>0\), the ball of the restricted metric is
\(\{y\in S:\rho(x,y)<r\}=S\cap\{y\in M:\rho(x,y)<r\}\). An open set of
\(S\) for the restricted metric is a union of such balls, so it has the
form \(S\cap U\) with \(U\) open in \(M\). Conversely let \(U\) be open in
\(M\) and \(x\in S\cap U\). Since \(U\) is a union of balls of \(M\), some
ball \(\{y:\rho(z,y)<r'\}\subseteq U\) contains \(x\); with
\(r=r'-\rho(z,x)>0\), the triangle inequality puts the ball
\(\{y:\rho(x,y)<r\}\) inside it. So \(S\cap U\) is the union, over its
points \(x\), of balls of the restricted metric centred at \(x\), and is
open in \(S\). Thus the open sets of \(S\) are exactly the sets
\(S\cap U\), \(U\) open in \(M\).

Write \(\mathcal B_\rho(S)\) for the sigma-algebra generated by these open
sets and \(\mathcal T=\{S\cap B:B\in\mathcal B(M)\}\). The family
\(\mathcal T\) is a sigma-algebra on \(S\): it contains \(S=S\cap M\);
\(S\setminus(S\cap B)=S\cap(M\setminus B)\); and
\(\bigcup_k(S\cap B_k)=S\cap\bigcup_kB_k\). It contains every open set of
\(S\), so \(\mathcal B_\rho(S)\subseteq\mathcal T\). In the other
direction, \(\mathcal G=\{B\subseteq M:S\cap B\in\mathcal B_\rho(S)\}\) is
a sigma-algebra on \(M\) by the same three identities, and it contains
every open \(U\subseteq M\), since \(S\cap U\) is open in \(S\). Hence
\(\mathcal B(M)\subseteq\mathcal G\), which says
\(\mathcal T\subseteq\mathcal B_\rho(S)\). Finally let
\(S\in\mathcal B(M)\). For \(B\in\mathcal B(M)\), \(S\cap B\) belongs to
\(\mathcal B(M)\) and is contained in \(S\); and every \(B\in\mathcal B(M)\)
with \(B\subseteq S\) equals \(S\cap B\in\mathcal T\). So
\(\mathcal T=\{B\in\mathcal B(M):B\subseteq S\}\).

(ii) By \cref{thm:p15-metric-borel-representation}, \(S\) with the
sigma-algebra \(\mathcal T\) has a Borel real coding \(\phi\) onto a
Borel set \(E\subseteq\mathbb R\), the empty map onto the empty set when
\(S=\varnothing\). By \cref{thm:p15-w02c-p15-004-nice-coding}(i),
\(\phi\) is a nice map for \((S,\mathcal T)\) with \(\phi(S)=E\) Borel. By
part~(i), \(\mathcal T\) is the source's \(\mathcal S\), so
\((S,\mathcal S)\) is nice.
\end{proof}
\end{annexitem}

\begin{annexitem}{annex:p15-w02c-p15-004-nice-extension}{thm:p15-w02c-p15-004-nice-extension}{Uniqueness, existence and a counterexample}
\begin{proof}
Throughout, measurability of a map into a product is checked on the
coordinate sets that generate the product sigma-algebra, by
\cref{thm:p15-measurability-criterion}. In particular, for \(m\geq n\)
the projection \(\prod_{i\leq m}S_i\to\prod_{i\leq n}S_i\) and the
projection \(\Pi_n:\prod_{i\geq1}S_i\to\prod_{i\leq n}S_i\) onto the first
\(n\) coordinates are measurable, since each pulls the coordinate set
\(\{y:y_i\in A\}\) back to the coordinate set with the same \(i\) and
\(A\).

(i) Let \(\mathcal C\) be the family of cylinders
\(\Pi_n^{-1}(E)=\{x:(x_1,\ldots,x_n)\in E\}\), \(n\geq1\), \(E\)
measurable in \(\prod_{i\leq n}S_i\). Each is measurable. Each coordinate
set \(\{x:x_i\in A\}\) is the cylinder with \(n=i\) and
\(E=\{y:y_i\in A\}\), so \(\sigma(\mathcal C)\) is the product
sigma-algebra. Given two cylinders with indices \(n\leq m\), the first
equals \(\Pi_m^{-1}(E')\), where \(E'\) is the inverse image of \(E\)
under the measurable projection \(\prod_{i\leq m}S_i\to\prod_{i\leq n}S_i\);
so their intersection is \(\Pi_m^{-1}(E'\cap F)\), again a cylinder, and
\(\mathcal C\) is a pi-system. Two probabilities with the same
finite-dimensional laws agree on \(\mathcal C\), hence on the product
sigma-algebra by \cref{thm:p15-measure-uniqueness}.

(ii) By \cref{thm:p15-w02c-p15-004-nice-coding}(ii) each nice map with
Borel image is a Borel real coding. The consistency assumed here is the
consistency under dropping the last coordinate, on all measurable
finite-product events, assumed in \cref{thm:p15-coded-space-extension},
which therefore gives a probability on \(\prod_{i\geq1}S_i\) with the laws
\(\nu_n\); part~(i) gives its uniqueness.

(iii) \emph{Rotations of \([0,1)\).} For real \(t\) put
\(\langle t\rangle=t-\lfloor t\rfloor\in[0,1)\), so that
\(\langle\langle a\rangle+b\rangle=\langle a+b\rangle\), the two
arguments differing by an integer. For \(r,s\in[0,1)\) put
\(r\oplus s=\langle r+s\rangle\) and \(r\ominus s=\langle r-s\rangle\),
and define \(\tau_r:[0,1)\to[0,1)\) by \(\tau_r(x)=\langle x+r\rangle\).
Then \(\tau_s\circ\tau_r=\tau_{r\oplus s}\) and \(\tau_0\) is the
identity, so \(\tau_r\) is a bijection with inverse
\(\tau_{\langle-r\rangle}\). For \(x\in[0,1-r)\), \(\tau_r(x)=x+r\), and for
\(x\in[1-r,1)\), \(\tau_r(x)=x+r-1\). Hence for a Borel set
\(B\subseteq[0,1)\),
\[
 \tau_r(B)=\bigl((B\cap[0,1-r))+r\bigr)\cup\bigl((B\cap[1-r,1))+(r-1)\bigr),
\]
a union of two translates of Borel sets, disjoint because \(\tau_r\) is
injective. By the translation invariance in
\cref{thm:p15-lebesgue-unit-interval}, \(\tau_r(B)\) is Borel and
\(\lambda(\tau_r(B))=\lambda(B\cap[0,1-r))+\lambda(B\cap[1-r,1))=\lambda(B)\).

\emph{A partition by rational rotations.} On \([0,1)\) let \(x\sim y\) when
\(x-y\in\mathbb Q\); this is an equivalence relation. Its classes are
nonempty, so the Axiom of Choice of \cref{sec:set-prerequisites-06} gives
a function choosing one member of each class; let \(V\) be the set of
chosen members. Put \(\mathbb Q_1=\mathbb Q\cap[0,1)\) and
\(V_q=\tau_q(V)\) for \(q\in\mathbb Q_1\). Three facts hold.
(a) The sets \(V_q\) cover \([0,1)\): for \(x\in[0,1)\) let \(v\in V\) be
the member chosen from the class of \(x\); then \(x-v\in\mathbb Q\), so
\(q=\langle x-v\rangle\in\mathbb Q_1\), and
\(\tau_q(v)=\langle v+\langle x-v\rangle\rangle=\langle x\rangle=x\).
(b) They are pairwise disjoint: if \(\tau_q(v)=\tau_{q'}(v')\) with
\(v,v'\in V\), then \((v+q)-(v'+q')\) is an integer, so \(v-v'\in\mathbb Q\),
\(v\) and \(v'\) were chosen from the same class, \(v=v'\), and \(q-q'\)
is an integer with \(q,q'\in[0,1)\), so \(q=q'\).
(c) \(\tau_s(V_q)=V_{q\oplus s}\) for \(q,s\in\mathbb Q_1\), since
\(\tau_s\circ\tau_q=\tau_{q\oplus s}\) and \(q\oplus s\in\mathbb Q_1\).

The set \(\mathbb Q_1\) is infinite, as it contains \(1/m\) for every
\(m\geq2\), and it is the image of a map from \(\mathbb N\): the rationals
are the image of the countable set
\(\mathbb Z\times(\mathbb N\setminus\{0\})\) under \((m,n)\mapsto m/n\), as in
the proof of \cref{thm:p15-w02c-p15-004-reals-continuum}, and values outside
\([0,1)\) may be replaced by \(0\). Listing the values of such a map in
order and deleting each value already listed gives a sequence
\(q_0,q_1,q_2,\ldots\) without repetition, which does not terminate because
\(\mathbb Q_1\) is infinite and in which every member of \(\mathbb Q_1\)
occurs. For \(n\geq1\) put
\[
 W_n=V_{q_0}\cup\cdots\cup V_{q_{n-1}},\qquad A_n=[0,1)\setminus W_n .
\]
Then \(A_1\supseteq A_2\supseteq\cdots\), and \(\bigcap_nA_n=\varnothing\):
by (a) each \(x\in[0,1)\) lies in some \(V_{q_k}\subseteq W_{k+1}\), so
\(x\notin A_{k+1}\).

\emph{Every Borel set containing \(A_n\) has full measure.} Fix \(n\geq1\)
and a Borel \(C\supseteq A_n\), and put \(K=[0,1)\setminus C\), a Borel
subset of \(W_n\). The set \(\Phi=\{q_i\ominus q_j:i,j<n\}\) is a finite
subset of \(\mathbb Q_1\), and \(0\ominus\eta\in\Phi\) whenever
\(\eta\in\Phi\). Choose \(s_1=0\) and, recursively,
\(s_{l+1}\in\mathbb Q_1\) outside the finite set
\(\{s_m\oplus\eta:m\leq l,\ \eta\in\Phi\}\). If \(m<l\) and
\(s_l\ominus s_m=\eta\in\Phi\), then \(s_l=s_m\oplus\eta\), which the
choice of \(s_l\) excludes; and \(s_m\ominus s_l=0\ominus(s_l\ominus s_m)\),
so it is not in \(\Phi\) either. Thus \(s_l\ominus s_m\notin\Phi\) whenever
\(l\neq m\). By (c), \(\tau_s(W_n)=\bigcup_{i<n}V_{q_i\oplus s}\). If
\(\tau_{s_l}(W_n)\) and \(\tau_{s_m}(W_n)\) met, (b) would give
\(q_i\oplus s_l=q_j\oplus s_m\) for some \(i,j<n\), hence
\(s_l\ominus s_m=q_j\ominus q_i\in\Phi\), which forces \(l=m\). So the
Borel sets \(\tau_{s_l}(K)\subseteq\tau_{s_l}(W_n)\), \(l\geq1\), are
pairwise disjoint subsets of \([0,1)\), each of measure \(\lambda(K)\).
For every \(L\),
\(L\,\lambda(K)=\lambda\bigl(\bigcup_{l\leq L}\tau_{s_l}(K)\bigr)\leq\lambda([0,1))=1\),
so \(\lambda(K)=0\) and \(\lambda(C\cap[0,1))=1-\lambda(K)=1\).

\emph{The sets \(A_n\) are not Borel.} For \(x\in[0,1)\), (a) gives
\(x\in V_q\) for some \(q\in\mathbb Q_1\), and with \(s=q\ominus q_0\),
(c) gives \(V_q=\tau_s(V_{q_0})\subseteq\tau_s(W_n)\). So the countably many
sets \(\tau_s(W_n)\), \(s\in\mathbb Q_1\), cover \([0,1)\). If \(A_n\) were
Borel, the previous paragraph with \(C=A_n\) would give
\(\lambda(A_n)=1\), so \(W_n=[0,1)\setminus A_n\) would be Borel of measure
zero, every \(\tau_s(W_n)\) would have measure zero, and by countable
subadditivity \(\lambda([0,1))=0\), which is false.

\emph{The factor spaces and their laws.} Let \(S_n=A_n\) with
\(\mathcal S_n=\{A_n\cap B:B\in\mathcal B(\mathbb R)\}\). By
\cref{thm:p15-w02c-p15-004-nice-coding}(iv) each inclusion is a nice map.
Define \(\lambda_n(A_n\cap B)=\lambda(B\cap[0,1))\). This is well defined:
if \(A_n\cap B=A_n\cap B'\), then \(A_n\) lies in the Borel set
\(\mathbb R\setminus(B\mathbin{\triangle}B')\), which therefore has measure
one in \([0,1)\), so \(\lambda((B\mathbin{\triangle}B')\cap[0,1))=0\) and
\(\lambda(B\cap[0,1))=\lambda(B'\cap[0,1))\). It is a probability measure:
\(\lambda_n(A_n)=\lambda([0,1))=1\); and if the sets \(A_n\cap B_k\) are
pairwise disjoint, the Borel sets
\(B'_k=B_k\setminus\bigcup_{j<k}B_j\) satisfy \(A_n\cap B'_k=A_n\cap B_k\)
and are pairwise disjoint, so
\(\lambda_n\bigl(\bigcup_k(A_n\cap B_k)\bigr)=\lambda\bigl(\bigcup_kB'_k\cap[0,1)\bigr)
=\sum_k\lambda(B'_k\cap[0,1))=\sum_k\lambda_n(A_n\cap B_k)\).

Since \(A_n\subseteq A_i\) for \(i\leq n\), the diagonal map
\(\Delta_n(u)=(u,\ldots,u)\) sends \(A_n\) into \(\prod_{i\leq n}S_i\). It
pulls the coordinate set \(\{y:y_i\in A_i\cap B\}\) back to
\(A_n\cap B\in\mathcal S_n\), so it is measurable, and
\(\nu_n(E)=\lambda_n(\Delta_n^{-1}(E))\) defines a probability law on
\(\prod_{i\leq n}S_i\). For measurable \(E\subseteq\prod_{i\leq n}S_i\),
write \(\Delta_n^{-1}(E)=A_n\cap B\) with \(B\) Borel. A point
\(u\in A_{n+1}\) has \(\Delta_{n+1}(u)\in E\times S_{n+1}\) exactly when
\(\Delta_n(u)\in E\), so
\(\Delta_{n+1}^{-1}(E\times S_{n+1})=A_{n+1}\cap A_n\cap B=A_{n+1}\cap B\) and
\[
 \nu_{n+1}(E\times S_{n+1})=\lambda_{n+1}(A_{n+1}\cap B)
   =\lambda(B\cap[0,1))=\lambda_n(A_n\cap B)=\nu_n(E).
\]
The sequence \((\nu_n)\) is consistent.

\emph{No extension.} Suppose \(P\) on \(\prod_{i\geq1}S_i\) had the
finite-dimensional laws \(\nu_n\). For \(n\geq2\) let
\(E_n=\{y\in\prod_{i\leq n}S_i:y_1=y_n\}\). Two distinct reals have a
rational strictly between them, so the complement of \(E_n\) is
\[
 \bigcup_{q\in\mathbb Q}\Bigl(\bigl(\{y:y_1<q\}\cap\{y:y_n>q\}\bigr)\cup
 \bigl(\{y:y_n<q\}\cap\{y:y_1>q\}\bigr)\Bigr),
\]
a countable union of intersections of coordinate sets such as
\(\{y:y_1\in A_1\cap(-\infty,q)\}\); hence \(E_n\) is measurable. Since
\(\Delta_n^{-1}(E_n)=A_n\), the set \(D_n=\{x:(x_1,\ldots,x_n)\in E_n\}\)
has \(P(D_n)=\nu_n(E_n)=\lambda_n(A_n)=1\). By countable subadditivity
applied to the complements, \(P\bigl(\bigcap_{n\geq2}D_n\bigr)=1\). But a
point \(x\) of \(\bigcap_{n\geq2}D_n\) would have \(x_1=x_n\in A_n\) for
every \(n\geq2\), and \(x_1\in A_1\), so \(x_1\in\bigcap_nA_n=\varnothing\).
The intersection is empty and cannot have probability one. No such \(P\)
exists.

\emph{Comparison with Borel real codings.} If every \((S_i,\mathcal S_i)\)
had a Borel real coding, then by
\cref{thm:p15-w02c-p15-004-nice-coding}(i) each would have a nice map with
Borel image, and part~(ii) would give an extension of \((\nu_n)\). So some
\((S_i,\mathcal S_i)\) has no Borel real coding, although it is nice
through its inclusion.
\end{proof}
\end{annexitem}

\begin{annexitem}{annex:p15-w02c-p15-004-reals}{thm:p15-w02c-p15-004-reals-continuum}{Ternary and rational-cut injections}
\begin{proof}
Define \(T:2^{\mathbb N}\to\mathbb R\) by
\(T(y)=\sum_{k\geq0}2y_k3^{-k-1}\); the series has nonnegative terms
bounded by those of \(\sum_k2\cdot3^{-k-1}=1\), so it converges. Let
\(y\neq y'\), let \(k\) be the least index with \(y_k\neq y'_k\), and
say \(y_k=1\), \(y'_k=0\). Since
\(\sum_{j>k}2\cdot3^{-j-1}=3^{-k-1}\),
\[
 T(y)-T(y')=2\cdot3^{-k-1}+\sum_{j>k}2(y_j-y'_j)3^{-j-1}
           \geq2\cdot3^{-k-1}-3^{-k-1}>0 .
\]
So \(T\) is injective.

The rationals are countable, being the image of the countable set
\(\mathbb Z\times(\mathbb N\setminus\{0\})\) under \((m,n)\mapsto m/n\);
fix an enumeration \(q_0,q_1,\ldots\) of \(\mathbb Q\). Define
\(D:\mathbb R\to2^{\mathbb N}\) by \(D(x)_j=1\) if \(q_j<x\) and
\(D(x)_j=0\) otherwise. If \(x<x'\), choose a positive integer \(n\) with
\(1/n<x'-x\) and let \(m\) be the least integer with \(m/n>x\); then
\(m/n\leq x+1/n<x'\). With \(m/n=q_j\), \(D(x)_j=0\) and
\(D(x')_j=1\). So \(D\) is injective. The injections \(D\) and \(T\)
in opposite directions, with the Cantor--Schr\"oder--Bernstein theorem
of \cref{sec:set-prerequisites-05}, give a bijection
\(\mathbb R\to2^{\mathbb N}\).

Finally, \(y\mapsto\{k:y_k=1\}\) is a bijection \(P\) of
\(2^{\mathbb N}\) with the power set of \(\mathbb N\). Suppose
\(\mathbb R\) were finite or countably infinite. It is nonempty, so there
would be a surjection \(r:\mathbb N\to\mathbb R\): a bijection
\(\mathbb N\to\mathbb R\) in the countably infinite case, and in the
finite case a bijection \(\{0,\ldots,n-1\}\to\mathbb R\) with \(n\geq1\),
extended by sending every \(j\geq n\) to its value at \(0\). Composing
\(r\), the bijection \(\mathbb R\to2^{\mathbb N}\) and \(P\) would give a
surjection of \(\mathbb N\) onto its power set, which Cantor's theorem of
\cref{sec:set-prerequisites-05} excludes. Hence \(\mathbb R\) is not
countable, and neither is \(2^{\mathbb N}\), which is in bijection with
it.
\end{proof}
\end{annexitem}

\begin{annexitem}{annex:p15-w02c-p15-004-borel-suslin}{thm:p15-w02c-p15-004-borel-suslin}{Closed sets, unions, intersections and the generated sigma-algebra}
\begin{proof}
\emph{Closed sets.} Let \(C\subseteq\mathbb R\) be closed and nonempty.
Fix a bijection \(e:\mathbb N\to\mathbb Z\). For \(k\geq1\) and
\(m\in\mathbb Z\) put \(J_{k,m}=[m2^{-k},(m+1)2^{-k}]\). For a finite
sequence \(s\) of length \(k\) put
\[
 C_s=C\cap J_{1,e(s_0)}\cap\cdots\cap J_{k,e(s_{k-1})},
 \qquad C_{()}=C,
\]
and let \(F=\{\alpha\in\mathcal N:C_{\alpha|k}\neq\varnothing
\hbox{ for every }k\}\). If \(\alpha\notin F\), some
\(C_{\alpha|k}\) is empty, and every \(\beta\in\mathcal N_{\alpha|k}\)
has the same empty \(C_{\beta|k}\); so \(F\) is closed.
For \(\alpha\in F\) the sets \(C_{\alpha|k}\), \(k\geq1\), are closed,
nonempty, decreasing, and each lies in an interval of length \(2^{-k}\).
Choosing \(c_k\in C_{\alpha|k}\), we have \(|c_j-c_k|\leq2^{-k}\) for
\(j\geq k\), so \((c_k)\) is Cauchy with a real limit \(x\). For each
\(k\) the terms \(c_j\), \(j\geq k\), lie in the closed set
\(C_{\alpha|k}\), so \(x\in C_{\alpha|k}\). Two points lying in every
\(C_{\alpha|k}\) are within \(2^{-k}\) of each other for every \(k\), so
\(x\) is the unique common point; put \(f(\alpha)=x\). If
\(\beta\in F\cap\mathcal N_{\alpha|k}\), then \(f(\alpha)\) and
\(f(\beta)\) both lie in \(C_{\alpha|k}\), so
\(|f(\beta)-f(\alpha)|\leq2^{-k}\); thus \(f\) is continuous, and
\(f[F]\subseteq C\). Conversely, for \(x\in C\) let
\(\alpha(k-1)=e^{-1}(\lfloor2^kx\rfloor)\) for \(k\geq1\). Then
\(x\in J_{k,\lfloor2^kx\rfloor}\) for every \(k\), so
\(x\in C_{\alpha|k}\) for every \(k\), \(\alpha\in F\), and
\(f(\alpha)=x\). Hence \(C=f[F]\). The empty set is a Suslin image by
definition.

\emph{Countable unions.} Let \(A=\bigcup_{n\in\mathbb N}A_n\) with each
\(A_n\) a Suslin image, and let \(I=\{n:A_n\neq\varnothing\}\). If \(I\)
is empty, \(A=\varnothing\). Otherwise write \(A_n=f_n[F_n]\) for
\(n\in I\). For \(\alpha\in\mathcal N\) let \(\alpha'(k)=\alpha(k+1)\), and
put
\[
 F=\{\alpha:\alpha(0)\in I,\ \alpha'\in F_{\alpha(0)}\},\qquad
 f(\alpha)=f_{\alpha(0)}(\alpha').
\]
If \(\alpha\notin F\), either \(\alpha(0)\notin I\), and then
\(\mathcal N_{\alpha|1}\cap F=\varnothing\), or
\(\alpha'\notin F_{\alpha(0)}\), and then some \(k\) has
\(\mathcal N_{\alpha'|k}\cap F_{\alpha(0)}=\varnothing\), whence
\(\mathcal N_{\alpha|k+1}\cap F=\varnothing\). So \(F\) is closed. For
\(\alpha\in F\) and \(\varepsilon>0\), take \(k\) from the continuity of
\(f_{\alpha(0)}\) at \(\alpha'\); every
\(\beta\in F\cap\mathcal N_{\alpha|k+1}\) has \(\beta(0)=\alpha(0)\) and
\(\beta'\in F_{\alpha(0)}\cap\mathcal N_{\alpha'|k}\), so
\(|f(\beta)-f(\alpha)|<\varepsilon\). Clearly
\(f[F]=\bigcup_{n\in I}A_n=A\).

\emph{Countable intersections.} Let \(A=\bigcap_{n\in\mathbb N}A_n\) with
each \(A_n\) a Suslin image. If some \(A_n\) is empty, so is \(A\).
Otherwise write \(A_n=f_n[F_n]\). Fix a bijection
\(\pi:\mathbb N\times\mathbb N\to\mathbb N\), and for
\(\alpha\in\mathcal N\) define \(\alpha_n(k)=\alpha(\pi(n,k))\). Given
\(n\) and \(k\), let \(m(n,k)\) exceed \(\pi(n,j)\) for all \(j<k\);
then \(\alpha|m(n,k)=\beta|m(n,k)\) implies
\(\alpha_n|k=\beta_n|k\). Put
\[
 F=\{\alpha:\alpha_n\in F_n\hbox{ and }f_n(\alpha_n)=f_0(\alpha_0)
     \hbox{ for every }n\},\qquad f(\alpha)=f_0(\alpha_0).
\]
Let \(\alpha\notin F\). If some \(\alpha_n\notin F_n\), take \(k\) with
\(\mathcal N_{\alpha_n|k}\cap F_n=\varnothing\); every
\(\beta\in\mathcal N_{\alpha|m(n,k)}\) has
\(\beta_n\in\mathcal N_{\alpha_n|k}\), so \(\beta_n\notin F_n\) and
\(\beta\notin F\). Otherwise \(f_n(\alpha_n)\neq f_0(\alpha_0)\) for some
\(n\); let \(\varepsilon\) be a third of the difference, take \(k\) so that
\(f_n\) differs from \(f_n(\alpha_n)\) by less than \(\varepsilon\) on
\(F_n\cap\mathcal N_{\alpha_n|k}\) and \(f_0\) differs from
\(f_0(\alpha_0)\) by less than \(\varepsilon\) on
\(F_0\cap\mathcal N_{\alpha_0|k}\), and let
\(m=\max\{m(n,k),m(0,k)\}\). For \(\beta\in\mathcal N_{\alpha|m}\),
either \(\beta_n\notin F_n\) or \(\beta_0\notin F_0\), and
\(\beta\notin F\); or both belong, and then
\(|f_n(\beta_n)-f_0(\beta_0)|>\varepsilon\), so again \(\beta\notin F\).
Thus \(F\) is closed. Continuity of \(f\) at \(\alpha\in F\) follows from
that of \(f_0\) at \(\alpha_0\), using \(m(0,k)\). Each
\(f(\alpha)=f_n(\alpha_n)\) lies in \(A_n\), so \(f[F]\subseteq A\).
If \(x\in A\), choose \(\alpha_n\in F_n\) with \(f_n(\alpha_n)=x\) and
define \(\alpha(\pi(n,k))=\alpha_n(k)\); this \(\alpha\) lies in \(F\)
and \(f(\alpha)=x\). Hence \(A=f[F]\); if \(A\) is empty it is a Suslin
image by definition.

\emph{Borel sets.} Let \(\mathcal D\) be the collection of sets
\(A\subseteq\mathbb R\) such that both \(A\) and \(\mathbb R\setminus A\)
are Suslin images. Let \(U\) be open. Its complement is closed, hence a
Suslin image. Each \(x\in U\) has some \((x-r,x+r)\subseteq U\), and then
rationals \(a<x<b\) in that interval give \(x\in[a,b]\subseteq U\); so
\(U\) is the union of the countably many closed intervals with rational
endpoints contained in it, and is a Suslin image by the first two steps.
Thus \(\mathcal D\) contains every open set, in particular \(\mathbb R\).
It is closed under complements by its symmetric definition. If
\(A_n\in\mathcal D\), then \(\bigcup_nA_n\) is a Suslin image by the
union step, and its complement \(\bigcap_n(\mathbb R\setminus A_n)\) is
one by the intersection step; so \(\bigcup_nA_n\in\mathcal D\). Hence
\(\mathcal D\) is a sigma-algebra containing the open sets, and it
contains \(\mathcal B(\mathbb R)\) by \cref{def:p15-borel-sets}.
\end{proof}
\end{annexitem}

\begin{annexitem}{annex:p15-w02c-p15-004-cond}{thm:p15-w02c-p15-004-condensation}{Removing countable pieces, splitting the rest}
\begin{proof}
The set of finite sequences of natural numbers is countable, as the
countable union over \(k\) of the countable sets \(\mathbb N^k\). Let
\(Z\) be the set of finite sequences \(s\) for which
\(f[F\cap\mathcal N_s]\) is countable, and put
\(F'=F\setminus\bigcup_{s\in Z}\mathcal N_s\). Then
\(f[F\setminus F']=\bigcup_{s\in Z}f[F\cap\mathcal N_s]\) is a countable
union of countable sets, hence countable, and so \(f[F']\) is
uncountable; in particular \(F'\neq\varnothing\).

\emph{Claim.} If \(F'\cap\mathcal N_s\neq\varnothing\), then
\(f[F'\cap\mathcal N_s]\) is uncountable. Indeed, a point
\(\alpha\in F'\cap\mathcal N_s\) lies in \(\mathcal N_s\) and in no
\(\mathcal N_t\) with \(t\in Z\), so \(s\notin Z\) and
\(f[F\cap\mathcal N_s]\) is uncountable; removing the countable set
\(f[(F\setminus F')\cap\mathcal N_s]\) leaves an uncountable subset of
\(f[F'\cap\mathcal N_s]\).

\emph{Splitting.} To each finite binary sequence \(u\) we assign a finite
sequence \(t_u\) of natural numbers with
\(F'\cap\mathcal N_{t_u}\neq\varnothing\), by recursion on the length of
\(u\). Put \(t_{()}=()\). Given \(t_u\), the claim provides
\(\alpha,\beta\in F'\cap\mathcal N_{t_u}\) with \(f(\alpha)\neq f(\beta)\).
Let \(r=|f(\alpha)-f(\beta)|/3\), and by continuity choose \(k>|t_u|\)
such that \(f\) maps \(F\cap\mathcal N_{\alpha|k}\) into
\(U_u=(f(\alpha)-r,f(\alpha)+r)\) and \(F\cap\mathcal N_{\beta|k}\) into
\(V_u=(f(\beta)-r,f(\beta)+r)\). The intervals \(U_u\) and \(V_u\) are
disjoint. Put \(t_{u0}=\alpha|k\) and \(t_{u1}=\beta|k\). Both extend
\(t_u\), are strictly longer, and meet \(F'\) at \(\alpha\) and \(\beta\)
respectively.

For \(y\in2^{\mathbb N}\), the sequences \(t_{y|j}\), \(j\geq0\), extend
one another with strictly increasing lengths, so there is a unique
\(\gamma_y\in\mathcal N\) extending all of them. If \(\gamma_y\notin F\),
some \(k\) has \(\mathcal N_{\gamma_y|k}\cap F=\varnothing\); but for
\(j\) with \(|t_{y|j}|\geq k\), the set
\(\mathcal N_{t_{y|j}}\subseteq\mathcal N_{\gamma_y|k}\) meets
\(F'\subseteq F\). Hence \(\gamma_y\in F\).

Let \(y\neq y'\), let \(u\) be their longest common initial segment, and
say \(y\) extends \(u0\) and \(y'\) extends \(u1\). Then
\(\gamma_y\in F\cap\mathcal N_{t_{u0}}\) and
\(\gamma_{y'}\in F\cap\mathcal N_{t_{u1}}\), so \(f(\gamma_y)\in U_u\)
and \(f(\gamma_{y'})\in V_u\), which are disjoint. Thus
\(y\mapsto f(\gamma_y)\) is an injection \(2^{\mathbb N}\to f[F]\).
\end{proof}
\end{annexitem}

\begin{annexitem}{annex:p15-w02c-p15-004-borel-cardinality}{thm:p15-w02c-p15-004-borel-cardinality}{Reduction to the line through the real coding}
\begin{proof}
By \cref{thm:p15-metric-borel-representation}, \(S\) has a Borel real
coding: a bijection \(\phi\) of \(S\) onto a Borel set
\(E\subseteq\mathbb R\) (\cref{def:p15-measurable-real-coding}). If \(S\)
is countable it is finite or countably infinite. Suppose \(S\) is
uncountable; then so is \(E\), and in particular \(E\neq\varnothing\).
By \cref{thm:p15-w02c-p15-004-borel-suslin}, \(E=f[F]\) for a closed
\(F\subseteq\mathcal N\) and a continuous \(f\), and
\cref{thm:p15-w02c-p15-004-condensation} gives an injection
\(2^{\mathbb N}\to E\). The inclusion \(E\to\mathbb R\) followed by the
bijection of \cref{thm:p15-w02c-p15-004-reals-continuum} is an injection
\(E\to2^{\mathbb N}\). By the Cantor--Schr\"oder--Bernstein theorem
of \cref{sec:set-prerequisites-05} there is a bijection
\(E\to2^{\mathbb N}\); composing \(\phi\), this bijection and the inverse
of the bijection \(\mathbb R\to2^{\mathbb N}\) gives a bijection
\(S\to\mathbb R\).

The three alternatives exclude one another: a finite set is not
countably infinite, and a set in bijection with \(\mathbb R\) is not
countable because \(\mathbb R\) is not, by
\cref{thm:p15-w02c-p15-004-reals-continuum}.
\end{proof}
\end{annexitem}

\begin{annexitem}{annex:p15-w02c-p15-004-birthday-error}{thm:p15-w02c-p15-004-birthday-error}{Bounds for the twenty-one factors}
\begin{proof}
Put \(x_j=j/365\) for \(1\leq j\leq21\); then \(0<x_j\leq21/365<1\).
For \(0<x<1\) the series \(-\log(1-x)=\sum_{k\geq1}x^k/k\) converges
and has positive terms. Dropping the terms with \(k\geq4\) gives
\(-\log(1-x)>x+x^2/2+x^3/3\). Since \(x^k/k<x^k/3\) for \(k\geq4\),
the geometric series gives
\(-\log(1-x)<x+x^2/2+\sum_{k\geq3}x^k/3=x+x^2/2+x^3/(3(1-x))\).

Let \(\Lambda=-\log p=\sum_{j=1}^{21}-\log(1-x_j)\) and
\(s_k=\sum_{j=1}^{21}x_j^k\). The power sums
\(\sum_{j=1}^{21}j=231\), \(\sum_{j=1}^{21}j^2=3311\) and
\(\sum_{j=1}^{21}j^3=231^2=53361\) give
\(s_1=231/365\), \(s_2=3311/365^2\) and \(s_3=53361/365^3\). Each
factor \(1-x_j\) is at least \(344/365\), so summing the two bounds
over \(j\) gives
\[
 s_1+\frac{s_2}2+\frac{s_3}3<\Lambda
   <s_1+\frac{s_2}2+\frac{365\,s_3}{3\cdot344}.
\]
In decimals, \(s_1=0.6328767\ldots\), \(s_2/2=0.0124263\ldots\),
\(s_3/3=0.0003657\ldots\) and \(365s_3/(3\cdot344)=0.0003881\ldots\),
so \(0.6456688<\Lambda<0.6456912\).

Put \(\kappa=484/730-\Lambda\). Since \(484/730=0.6630136\ldots\),
\(0.017322<\kappa<0.017345\). Then
\(a/p=e^{-484/730}/e^{-\Lambda}=e^{-\kappa}<1\), so \(a<p\), and
\((p-a)/p=1-e^{-\kappa}\).

For \(t>0\), \(t-t^2/2<1-e^{-t}<t\). Indeed, \(g(t)=e^{-t}-1+t\) has
\(g(0)=0\) and \(g'(t)=1-e^{-t}>0\), so \(e^{-t}>1-t\); and
\(h(t)=1-t+t^2/2-e^{-t}\) has \(h(0)=0\) and
\(h'(t)=-1+t+e^{-t}=g(t)>0\), so \(e^{-t}<1-t+t^2/2\). The function
\(t\mapsto t-t^2/2\) increases on \((0,1)\). Hence
\[
 \frac{p-a}p<\kappa<0.017345<0.0174
\]
and
\[
 \frac{p-a}p>0.017322-\frac{0.017322^2}2>0.017171>0.0171.
\]

For the source's value \(a'=e^{-0.6630}\): \(730\cdot0.6630=483.99\),
so \(484/730-0.6630=1/73000<0.0000137\). Thus
\(a'/p=e^{-\kappa'}\) with \(\kappa'=\kappa-1/73000\), and
\(0.0173083<\kappa'<\kappa<0.017345\). As before, \(a'<p\),
\((p-a')/p<\kappa'<0.0174\), and
\((p-a')/p>0.0173083-0.0173083^2/2>0.017158>0.0171\).
\end{proof}
\end{annexitem}

\begin{annexitem}{annex:p15-w02c-p15-004-dens}{thm:p15-w02c-p15-004-density-reading}{Singletons and a return probability}
\begin{proof}
(i) By \cref{thm:p15-weak-central-order-statistic}, \(Y_n\) has the
density \(f_{Y_n}\), so
\(\mathbb P(Y_n\in B)=\int_Bf_{Y_n}(t)\,dt\) for every Borel
\(B\subseteq\mathbb R\). The singleton \(\{y\}\) lies in
\([y-\eta,y+\eta]\) for every \(\eta>0\), so its Lebesgue measure is at
most \(2\eta\) for every \(\eta>0\) and hence zero. The integral of the
nonnegative function \(f_{Y_n}\) over a Lebesgue-null set is zero, so
\(\mathbb P(Y_n=y)=0\).

(ii) The point \(y=0\) satisfies \(|0|<\sqrt{2n}\), so the formula of
\cref{thm:p15-weak-central-order-statistic} applies and gives
\(f_{Y_n}(0)=C_n\cdot1\cdot\frac{2n+1}{2n}\sqrt{n/2}\). Each factor is
positive, \(C_n=4^{-n}\binom{2n}{n}\) included.

(iii) \(S_{2n}=0\) exactly when exactly \(n\) of the \(2n\) signs equal
\(1\). For a fixed set \(J\subseteq\{1,\ldots,2n\}\) with \(n\)
elements, the event that \(\varepsilon_i=1\) for \(i\in J\) and
\(\varepsilon_i=-1\) for \(i\notin J\) has probability \(2^{-2n}\) by
independence. These events are disjoint for distinct \(J\), and there
are \(\binom{2n}{n}\) such sets, so
\(\mathbb P(S_{2n}=0)=\binom{2n}{n}2^{-2n}=4^{-n}\binom{2n}{n}=C_n\).
\end{proof}
\end{annexitem}
